\section{Method}
\label{sec:methods}

\subsection{Overview and notation}
\label{sec:overview}

We infer the stellar mass $M$ as a function of absolute magnitude
$x\equiv M_G$ and metallicity $Z\equiv\mh$.  Binary systems are indexed by
$j$, their components by $k\in\{1,2\}$, and the nodes of a metallicity grid
by $q$.  Hats denote observed quantities; $\log_{10}$ is used for stellar-mass
logarithms and $\log$ for natural logarithms.

The analysis is modular \citep{LiuBayarriBerger2009Modularization,
Plummer2015Cuts} (Fig.~\ref{fig:workflow}).  Stage~1a calibrates the relative,
magnitude-dependent behaviour of the XP metallicities from within-binary
differences.  Stage~1b combines the calibrated metallicities with the PARSEC
colour--magnitude relation and a population distribution to obtain, for each
system, a posterior probability mass $\overline P_{jq}$ for its shared
metallicity on a fixed grid $\{Z_q\}$.  Stage~2 predicts the two component
masses at every grid node, evaluates the dynamical likelihood of $u_j$, and
sums over nodes with weights $\overline P_{jq}$.  Information flows in one
direction only: the dynamics do not recalibrate the metallicity zero point,
the colour scatter, or the population distribution.

\begin{figure*}[t]
\centering
\begin{tikzpicture}[
  box/.style={draw,rounded corners,align=center,minimum height=11mm,
              text width=34mm,fill=blue!5,font=\small},
  data/.style={box,fill=gray!10},
  model/.style={box,fill=orange!10},
  result/.style={box,fill=green!10},
  end/.style={box,fill=red!15},
  arrow/.style={-{Latex[length=2.2mm]},thick}
]
\node[data]   (spec)   at (0, 2.6) {XP metallicities\\$\widehat z_{j1},\widehat z_{j2}$, $\sigma_{z,jk}$};
\node[data]   (cmd)    at (0, 1.0) {\gaia\ CMD\\$M_G,(G_{\rm BP}-G_{\rm RP})_0$};
\node[model]  (stage1) at (4.1, 1.8) {Stage 1a: relative XP calibration\\Stage 1b: shared $Z_j$ + CMD + population};
\node[result] (weights)at (8.2, 1.8) {System metallicity posterior\\$\overline P_{jq}$};
\node[data]   (phot)   at (0,-0.8) {Component magnitudes\\$x_{j1},x_{j2}$};
\node[model]  (mlr)    at (4.1,-0.8) {PARSEC-relative\\hard-monotone MLR\\$g(x,Z)$};
\node[result] (mass)   at (8.2,-0.8) {Total mass and scale\\$M_{{\rm tot},jq}$, $s_{jq}$};
\node[data]   (dyn)    at (0,-2.6) {Wide-binary dynamics\\$u_j,\sigma_{u,j}$};
\node[model]  (rice)   at (4.1,-2.6) {Rice-convolved\\normal/outlier likelihood\\shape $(B,u_c,C)$};
\node[end]    (post)   at (4.1,-4.5) {Stage-2 posterior\\MLR, shape, $f_{\rm out}$};
\draw[arrow] (spec) -- (stage1);
\draw[arrow] (cmd) -- (stage1);
\draw[arrow] (stage1) -- (weights);
\draw[arrow] (phot) -- (mlr);
\draw[arrow] (mlr) -- (mass);
\draw[arrow] (mass) -- (rice);
\draw[arrow] (dyn) -- (rice);
\draw[arrow] (rice) -- (post);
\draw[arrow] (weights.east) -- ++(7mm,0) |- node[pos=0.42,right,align=left]
  {metallicity\\marginalisation} (rice.east);
\end{tikzpicture}
\caption{Structure of the inference.  Stage~1 returns a posterior probability
mass over the shared metallicity of every system.  Stage~2 propagates this
uncertainty into the dynamical fit of the MLR but does not feed back into
Stage~1.}
\label{fig:workflow}
\end{figure*}

\subsection{Stage 1: system metallicities}
\label{sec:stage1}

\subsubsection{XP measurement model}

Each star carries an independent Student-$t$ measurement of the shared system
metallicity $Z_j$,
\begin{equation}
 \widehat z_{jk}\mid Z_j
 \sim \StudentT_{\nu}\!\left[
 Z_j+\delta_z+b_G(x_{jk}),\
 \left(\sigma_{z,jk}^2+s_z^2\right)^{1/2}\right]
 \equiv\mathcal{L}^{\mathrm{XP}}_{jk}(Z_j),
 \label{eq:spec-model}
\end{equation}
where the second and third arguments are the location and scale.  The
function $b_G(x)$ is a magnitude-dependent relative bias, $s_z$ an additional
per-star scatter, $\nu$ the degrees of freedom, and $\delta_z$ a global zero
point.  In Stage~1a, $b_G$, $s_z$, and $\nu$ are estimated by maximum
likelihood from the within-binary differences $\widehat z_{j1}-\widehat z_{j2}$,
which do not depend on $Z_j$ or $\delta_z$; they are then held fixed.

\subsubsection{Colour--magnitude likelihood}

The zero point $\delta_z$ is not constrained by metallicity differences.  It
is set in Stage~1b by requiring the observed colours to agree with the PARSEC
colour $C_{\rm P}(x,Z)$ at the inferred metallicity.  For a component with
formal colour uncertainty
$\sigma_{C,jk}=(2.5/\ln10)\,(\mathrm{SNR}_{\rm BP}^{-2}+
\mathrm{SNR}_{\rm RP}^{-2})^{1/2}$,
\begin{equation}
 C_{jk}\mid Z_j\sim\StudentT_{4}\!\left[
 C_{\rm P}(x_{jk},Z_j),\ \left(\sigma_{C,jk}^2+s_C^2\right)^{1/2}\right]
 \equiv\mathcal{L}^{\rm CMD}_{jk}(Z_j),
 \label{eq:cmd}
\end{equation}
with an excess scatter $s_C\sim\text{Half-Normal}(0.1~{\rm mag})$.  Because
PARSEC colours of cool dwarfs are less reliable, only components whose PARSEC
effective temperature exceeds 4000~K at every grid metallicity enter this term:
\begin{equation}
 m_{jk}=\indicator\!\left[
 \min_{q}T_{\rm eff,P}(x_{jk},Z_q)\ge4000\ {\rm K}\right].
 \label{eq:warm-mask}
\end{equation}
The mask is independent of the proposed metallicity, so the set of stars that
enter the CMD term does not change during sampling.  It retains 12\,167 stars;
8\,247 systems have at least one and 3\,920 have two such components.  The
PARSEC reference surface is the average over isochrones with
$9.3<\log_{10}(\mathrm{age/yr})<10.0$, so age is not a latent variable.

\subsubsection{System posteriors}

The population distribution $p_{\rm pop}(Z\mid\bm\omega)$ is a six-component
mixture with Dirichlet$(1,\ldots,1)$ weights $\bm\omega$, and
$\delta_z\sim\Normal(0,0.3~{\rm dex})$.  On an 81-node grid spanning
$-1.0\le Z\le0.6$, the posterior probability mass for system $j$ is
\begin{equation}
 \overline P_{jq}\propto w_q\,
 \Big\langle p_{\rm pop}(Z_q\mid\bm\omega)
 \prod_{k=1}^{2}\mathcal{L}^{\rm XP}_{jk}(Z_q)
 \left[\mathcal{L}^{\rm CMD}_{jk}(Z_q)\right]^{m_{jk}}
 \Big\rangle_{\bm\psi},
 \label{eq:zpost}
\end{equation}
where $w_q$ are trapezoidal weights, $\sum_q\overline P_{jq}=1$, and the
average is over the Stage-1b posterior of $\bm\psi=(\delta_z,s_C,\bm\omega)$.
The global parameters are constrained by the product of the per-system
evidences.  Stage~1b was sampled with the No-U-Turn Sampler
\citep[NUTS;][]{HoffmanGelman2014NUTS} \unsure{(chains and draws of the
formal run)}.

\subsection{Stage 2: a PARSEC-relative hard-monotone MLR}
\label{sec:mlr}

\subsubsection{Spline representation}

Let $g_{\rm P}(x,Z)=\log_{10}[M_{\rm P}(x,Z)/\msun]$ be the PARSEC mass
surface.  The fitted surface is a tensor-product cubic B-spline
\citep{deBoor2001Splines},
\begin{equation}
 g(x,Z)=\sum_{i=1}^{K_x}\sum_{h=1}^{K_Z}\Theta_{ih}B_i(x)C_h(Z)
       =g_{\rm P}^{*}(x,Z)+\Delta_*(x,Z),
 \label{eq:g-total}
\end{equation}
where $g_{\rm P}^*=\sum\Pi_{ih}B_iC_h$ is the least-squares projection of
$g_{\rm P}$ onto the same basis and $\Delta_*=\sum D_{ih}B_iC_h$ with
$D=\Theta-\Pi$ is the fitted correction.  The projection error
$g_{\rm P}^*-g_{\rm P}$ has an rms of 0.002~dex and a maximum of 0.019~dex.
We use $K_x\times K_Z=8\times4$ basis functions with clamped knot vectors
$\bm\kappa_x=(3.5^{(4)},5.5,7.5,9.5,11.5,13.5^{(4)})$~mag and
$\bm\kappa_Z=(-1.0^{(4)},0.6^{(4)})$~dex, where superscripts give
multiplicities.  The metallicity dependence is therefore a single cubic
polynomial segment.

\subsubsection{Hard monotonicity}

At fixed metallicity a fainter star must not be more massive, and at fixed
magnitude we require that a more metal-rich star is not less massive.  We
impose
\begin{equation}
 \Theta_{i+1,h}\le\Theta_{ih},\qquad
 \Theta_{i,h+1}\ge\Theta_{ih},
 \label{eq:coefficient-order}
\end{equation}
which, because the derivatives of a B-spline are non-negative combinations of
adjacent coefficient differences, guarantees $\partial g/\partial x\le0$ and
$\partial g/\partial Z\ge0$ everywhere in the domain.  The constraints act on
the total surface $\Theta$, not on the correction $D$, which may take either
sign.  Appendix~\ref{app:monotone} gives the constructive parametrisation.

\subsubsection{Regularisation and solar anchor}

The correction is shrunk towards PARSEC and towards smoothness through
\begin{equation}
\begin{aligned}
 \log\pi(D)={}&-\frac{1}{2\tau_D^2}\sum_{i,h}D_{ih}^2
 -\frac{1}{2}\sum\left(\frac{D^{(2,x)}_{ih}}{\tau_x}\right)^2\\
 &-\frac{1}{2}\sum\left(\frac{D^{(2,Z)}_{ih}}{\tau_Z}\right)^2
 -\frac{1}{2}\sum\left(\frac{D^{(xZ)}_{ih}}{\tau_{xZ}}\right)^2+{\rm const},
\end{aligned}
 \label{eq:shrink}
\end{equation}
where $D^{(2,x)}$ and $D^{(2,Z)}$ are second differences and $D^{(xZ)}$ the
mixed difference of the coefficients, normalised by the Greville-abscissa
spacing so that they retain units of dex \citep{EilersMarx1996PSplines}.  We
adopt $\tau_D=0.10$~dex and $\tau_x=\tau_Z=\tau_{xZ}=0.05$~dex.  The absolute
mass scale is anchored softly at the Sun,
\begin{equation}
 g(4.67,0)\sim\Normal\!\left(0,\ [0.01/\ln10]^2\right),
 \label{eq:solar}
\end{equation}
i.e. a 1\% constraint on the mass of a solar-luminosity, solar-metallicity
star.

\subsection{Dynamical likelihood}
\label{sec:dynamics}

\subsubsection{Normal and outlier components}

For a bound binary on a Keplerian orbit, the projected relative velocity and
separation combine to
$u=\sqrt{M_{\rm tot}/\msun}\;\tilde u$, where $\tilde u$ depends only on
eccentricity, orbital phase, and orientation.  Marginalising these over an
assumed orbital population gives a distribution $p(\tilde u)$ that is common
to all systems.  We parametrise it with the functional form of
\citet{Hwang2024DynamicalMasses},
\begin{equation}
 p_{\rm N}(\tilde u\mid B,u_c,C)\propto\tilde u\,
 \exp\!\left[-B\tilde u^2-\exp\!\left(\frac{\tilde u-u_c}{C}\right)\right],
 \qquad 0<\tilde u\le80,
 \label{eq:good}
\end{equation}
normalised on its support.  The reference values
$(B,u_c,C)=(2.544\times10^{-3},35.67,3.1)$ were obtained by fitting this form
to simulated $\tilde u$ for an orbital population
\unsure{(specify eccentricity distribution and separation range used in the
simulation)}.

Pairs that are not simple bound binaries, such as hierarchical triples with an
unresolved inner companion or residual chance alignments, produce velocities
that do not scale with the photometric mass.  We describe them with an outlier
component defined directly on the latent observed-frame amplitude $v$,
\begin{equation}
 p_{\rm O}(v)=\frac{\phi[(v-40)/13]}
 {13\,[2\Phi(40/13)-1]}\,
 \indicator(0\le v\le80),
 \label{eq:out}
\end{equation}
a normal distribution with mean 40 and dispersion $13~\kmsau$ truncated to
$[0,80]$, independent of the mass and metallicity of the system.

\subsubsection{Measurement model}

The measured $u_j$ is the modulus of a two-dimensional velocity difference
with approximately isotropic Gaussian errors, so its conditional density
given the true amplitude $v$ is the Rice distribution \citep{Rice1945},
\begin{equation}
 R(u\mid v,\sigma)=\frac{u}{\sigma^2}
 \exp\!\left[-\frac{u^2+v^2}{2\sigma^2}\right]
 I_0\!\left(\frac{uv}{\sigma^2}\right).
 \label{eq:rice}
\end{equation}
For a mass scale $s=\sqrt{M_{\rm tot}/\msun}$ the two component likelihoods
are
\begin{align}
 L_{{\rm N},j}(s)&=\int_0^{80s}R(u_j\mid v,\sigma_{u,j})\,
 p_{\rm N}\!\left(\frac{v}{s}\right)\frac{\dd v}{s},
 \label{eq:LN}\\
 L_{{\rm O},j}&=\int_0^{80}R(u_j\mid v,\sigma_{u,j})\,p_{\rm O}(v)\dd v,
 \label{eq:LO}
\end{align}
and the mixture is
$L_j(s)=(1-f_{\rm out})L_{{\rm N},j}(s)+f_{\rm out}L_{{\rm O},j}$ with
$f_{\rm out}\sim{\rm Beta}(3,12)$.

\subsubsection{Metallicity marginalisation}

At metallicity node $Z_q$ the component masses and mass scale are
$M_{jkq}=\msun10^{g(x_{jk},Z_q)}$ and
$s_{jq}=[(M_{j1q}+M_{j2q})/\msun]^{1/2}$.  The Stage-2 likelihood of system
$j$ is
\begin{equation}
 \mathcal L_j^{\rm dyn}=\sum_{q=1}^{81}\overline P_{jq}\,L_j(s_{jq}),
 \label{eq:dyn-final}
\end{equation}
which propagates the metallicity uncertainty of each system into its mass.
Neither the population distribution nor the quadrature weights are applied a
second time, because both are contained in $\overline P_{jq}$.

\subsection{Metallicity bins and velocity-distribution shape}
\label{sec:bins}

The shape parameters $(B,u_c,C)$ strongly affect the inferred mass scale
(Sect.~\ref{sec:fixedshape}), and they need not be the same for binaries of
different metallicity, since the wide-binary fraction and separation
distribution depend on metallicity \citep{Hwang2021}.  We therefore fit four
metallicity bins independently, each with its own MLR surface, shape
parameters, and outlier fraction.  Systems are assigned by the median of
their metallicity posterior to the bins $[-1.0,-0.5)$, $[-0.5,0.0)$,
$[0.0,0.3)$, and $[0.3,0.6]$, which contain 186, 9\,618, 4\,687, and 385
systems.  Within a bin every system retains its full 81-node posterior in
Eq.~(\ref{eq:dyn-final}).  The MLR of each bin is reported at the bin
midpoint.

The shape parameters receive uniform priors in $\log B$, $\log u_c$, and
$\log C$ within the box $B\in[0.7,3.3]\times10^{-3}$, $u_c\in[24,44]$, and
$C\in[2.5,13.5]$.  To make sampling tractable, Eq.~(\ref{eq:LN}) is
precomputed for every system on a logarithmic grid of mass scales
$s\in[0.25,2.5]$ at $3\times3\times3$ nodes spanning this box.  At a general
point, the normal-component likelihood is the probability mixture of the 27
node likelihoods with trilinear weights in the logarithmic shape coordinates.
This mixture is itself a normalised distribution in $u$, although it is not
identical to Eq.~(\ref{eq:good}) evaluated at the interpolated parameters.

\subsection{Priors and posterior sampling}
\label{sec:inference}

The MLR is parametrised by 35 unconstrained coordinates
(Appendix~\ref{app:monotone}), to which the three shape parameters are added
(Table~\ref{tab:priors}).  Each bin was sampled with NUTS
\citep{HoffmanGelman2014NUTS} using four chains of 1\,500 warm-up and 2\,000
retained draws.  Each bin was sampled twice with independent seeds and
jittered initial values.  All Stage-2 runs completed without divergent
transitions.  Convergence diagnostics are listed in
Table~\ref{tab:diagnostics}.

\begin{table}
\caption{Stage-2 parameters and priors.}
\label{tab:priors}
\centering
\small
\begin{tabular}{@{}ll@{}}
\toprule
Parameter & Prior \\
\midrule
$c_0=\Theta_{11}$ & $\Normal(\Pi_{11},0.10)$ dex \\
$a_h,\ b_i,\ r_{ih}$ & $\Normal(0,1)$ \\
$\log\lambda_x$ & $\Normal(\log0.186,0.7)$ \\
$\log\lambda_Z$ & $\Normal(\log0.131,0.7)$ \\
$\log B$ & $\mathcal U(\log0.0007,\log0.0033)$ \\
$\log u_c$ & $\mathcal U(\log24,\log44)$ \\
$\log C$ & $\mathcal U(\log2.5,\log13.5)$ \\
$f_{\rm out}$ & Beta$(3,12)$ \\
\midrule
$\tau_D$; $\tau_x,\tau_Z,\tau_{xZ}$ & 0.10; 0.05 dex (fixed) \\
Solar anchor & $g(4.67,0)=0\pm0.0043$ dex \\
\bottomrule
\end{tabular}
\tablefoot{$\lambda_x$ and $\lambda_Z$ set the scale of the monotone
increments (Appendix~\ref{app:monotone}); their reference values are the mean
PARSEC coefficient increments.}
\end{table}
